% ----------------------------------------------------------------------
\section{Integration with the BX3 Framework}
\label{sec:rs-integration}
% ----------------------------------------------------------------------

Recursive Spawning is not a standalone delegation mechanism. Its correctness
guarantees are derived entirely from the structural properties enforced by the
BX3 three-layer architecture operating in mandatory concert.  This section
specifies the three integration points and establishes the unified invariant
from which all downstream correctness proofs descend.

% -----------------------------------------------------------------------
\subsection{Purpose Layer as Accountability Anchor}
\label{sec:rs-purpose-anchor}

Every node $n$ in a BX3 spawn tree carries a human anchor $h(n)$ established at
node provisioning and structurally non-modifiable for the node's entire
operational lifetime.  The Purpose layer is the architectural locus of this
anchor: it encodes the principal's intent, defines the operating range $R(n)$,
and serves as the final machine layer before the human principal in any
escalation chain.

Three integration points connect Recursive Spawning to the Purpose layer:

\begin{enumerate}
  \item[(a)] \textbf{Anchor inheritance at spawn.}  When parent node $p$ spawns
    child $c$, the spawn directive encodes $h(c) = h(p)$.  The Fact layer
    pre-commit hook enforces this as a structural constraint before the child
    node is provisioned.  The consequence is that \emph{every node in the
    spawn tree, regardless of depth, points to the same human principal} as
    its ultimate accountability terminus.

  \item[(b)] \textbf{Escalation terminus.}  The Bailout Protocol traverses the
    parent-pointer chain $\alpha(n), \alpha(\alpha(n)), \ldots$ until reaching
    $h(n)$.  Because $h(n)$ is always a human principal and never a machine
    actor, the escalation path never terminates at an intermediate machine
    node.  The Purpose layer at each node is the last machine layer before the
    human anchor --- it holds the principal's intent, not a delegated subset
    of it.

  \item[(c)] \textbf{Scope adjudication.}  When the Bounds Engine detects a
    delegation scope inheritance violation, the exception is escalated to the
    Purpose layer of the parent node for adjudication.  The Purpose layer alone
    has authority to redefine $R(p)$ or to authorize an exception to the
    inheritance constraint on behalf of the human principal.  No machine actor
    may grant this authorization.
\end{enumerate}

The combined effect of (a)--(c) is that the Purpose layer serves as a double
anchor: it anchors the spawn tree's accountability chain at the top (human
principal $\rightarrow$ root node's Purpose layer), and it anchors every
escalation path at its terminus.  The spawn tree is not a tree of autonomous
agents --- it is a tree of accountability structures, each anchored to the
same human principal.

% -----------------------------------------------------------------------
\subsection{Bounds Engine: Depth Enforcement}
\label{sec:rs-bounds-depth}

The Bounds Engine enforces the Hierarchy Finiteness Invariant as a deterministic
pre-commit check on every spawn event.  When a node in state \texttt{CAN\_SPAWN}
receives an $e_{\text{spawn}}$ event, it transitions to \texttt{SPAWN\_PENDING}
and issues a depth-confirmation request to the Bounds Engine.  The Bounds
Engine evaluates $\text{depth-ok}(c, p)$ by computing $d(c) = d(p) + 1$ and
comparing against $D_{\max}$.

\begin{itemize}
  \item \textbf{Pass:} Bounds Engine emits $e_{\text{depth-ok}}$; spawn
    proceeds.
  \item \textbf{Fail:} Bounds Engine emits $e_{\text{depth-violation}}$,
    transitions the spawning state machine to \texttt{SPAWN\_LOCKED}, triggers
    the Bailout Protocol, and records a \texttt{depth-rejected} entry in the
    Forensic Ledger.
\end{itemize}

The depth check is a deterministic integer comparison.  Either $d(c) \leq
D_{\max}$ or it does not; the Bounds Engine enforces the condition without
discretion.  Critically, $D_{\max}$ is a deployment parameter established at
system initialization and is not negotiable at runtime.  No machine actor,
including the Bounds Engine itself, may raise $D_{\max}$ to accommodate a
non-compliant spawn.

The delegation scope inheritance check is enforced by the Bounds Engine as a
deterministic set comparison: $R(c) \subset R(p)$.  If the subset relation
fails, the Bounds Engine emits $e_{\text{scope-violation}}$, transitions the
spawning state machine to \texttt{SPAWN\_LOCKED}, and triggers the Bailout
Protocol.  Scope and depth are evaluated in sequence --- scope first, then
depth --- ensuring the first detected violation is the one recorded in the
Forensic Ledger, which matters for accountability attribution during human
review.

% -----------------------------------------------------------------------
\subsection{Fact Layer: Forensic Ledger}
\label{sec:rs-fact-fledger}

The Fact layer provides two structural guarantees for Recursive Spawning:
immutable parent-pointer enforcement and complete, cryptographically chained
event recording.

\medskip\noindent\textbf{Immutable parent pointer.}  The parent-pointer field
$\alpha(v)$ is a read-only property of the node's Fact layer worksheet.  Any
attempt to write to $\alpha(v)$ at runtime is intercepted by the pre-commit hook
and rejected.  The rejection is recorded as a \texttt{parent-pointer-write-denied}
entry, and the Bailout Protocol is triggered if the attempted write was
initiated by a machine actor.  This enforcement ensures the accountability chain
is never redirected after node creation.

\medskip\noindent\textbf{Cryptographically chained forensic record.}  Every
spawn event, every state transition, every depth-confirmation result, every
scope-inheritance evaluation, and every human directive issued in response to a
spawn-related exception is recorded as a Forensic Ledger entry.  Each entry $f$
carries the SHA-256 hash of the preceding entry, forming an append-only chain.
Any modification, deletion, or insertion of an entry breaks the hash chain and
is detectable by any auditor without access to the ledger's contents.

The Forensic Ledger entries relevant to Recursive Spawning include:
\begin{itemize}\itemsep=0pt
  \item \texttt{spawn-proposed}$(p, c, R(c), d(c))$,
  \item \texttt{depth-confirmed}$(c, d(c))$ / \texttt{depth-rejected}$(c, d(c), D_{\max})$,
  \item \texttt{scope-confirmed}$(c, R(c) \subset R(p))$ / \texttt{scope-rejected}$(c, R(c), R(p))$,
  \item \texttt{child-provisioned}$(p, c, \alpha(c)=p, h(c)=h(p))$,
  \item \texttt{parent-pointer-write-denied}$(v, \alpha_{\text{old}}, \alpha_{\text{new}})$,
  \item \texttt{spawn-locked}$(v, \text{violation-type})$.
\end{itemize}

The Forensic Ledger is not a conventional audit log that records actions after
they are taken --- it is an enforcement mechanism that queries current state
before subsequent actions are permitted.  The Fact layer pre-commit hook
verifies the current state of the spawning state machine before committing any
spawn event, ensuring the ledger is the authoritative source of truth for the
runtime state of the spawn tree.

% -----------------------------------------------------------------------
\subsection{Unified Integration Invariant}
\label{sec:rs-unified-inv}

The three-layer integration yields a single structural invariant that is the
foundation of all downstream correctness guarantees:

\begin{invariant}[Spawn Tree Accountability Invariant]
\label{inv:spawn-accountability}
For every node $v$ in a BX3 spawn tree, the accountability chain
$\alpha^{(k)}(v) = h(v)$ with $k \leq D_{\max}$ is:
\begin{enumerate}
  \item[(i)] \textbf{Immutable} --- $\alpha(v)$ is never modified after node
    creation;
  \item[(ii)] \textbf{Complete} --- every spawn event and state transition is
    recorded in the Forensic Ledger;
  \item[(iii)] \textbf{Terminating} --- the chain terminates at a human
    principal $h(v)$, never at a machine actor;
  \item[(iv)] \textbf{Enforceable} --- any violation of (i), (ii), or (iii)
    triggers the Bailout Protocol before the system state advances.
\end{enumerate}
\end{invariant}

This invariant is enforced by the pre-commit hooks of all three layers operating
in concert.  The Purpose layer provides the human anchor; the Bounds Engine
provides depth and scope enforcement; the Fact layer provides the immutable
parent pointer and the cryptographically chained forensic record.  Together they
enforce Invariant~\ref{inv:spawn-accountability}, and it is this invariant that
enables the correctness properties established in the next section.
